% Bagian: first-order-logic
% Bab: syntax-and-semantics
% Subbagian: introduction

\documentclass[../../../include/open-logic-section]{subfiles}

\begin{document}

\olfileid[id]{fol}{syn}{int}

\olsection{Pendahuluan}

Untuk mengembangkan teori dan metateori logika orde pertama, mula-mula kita
harus mendefinisikan sintaks dan semantik ungkapan-ungkapannya. Ungkapan logika
orde pertama terdiri atas suku dan !!{formula}. Suku dibentuk dari
!!{variable}, !!{constant}, dan !!{function}. Selanjutnya, !!^{formula}
dibentuk dari !!{predicate} bersama suku (yang membentuk !!{formula} terkecil,
atau ``atomik''), lalu dari !!{formula} atomik kita dapat membentuk formula
yang lebih kompleks dengan menggunakan penghubung logis dan kuantor. Ada
banyak cara berbeda untuk menetapkan kaidah pembentukan; di sini kita hanya
memberikan salah satu cara. Sistem lain akan memilih simbol yang berbeda,
memilih himpunan penghubung primitif yang berbeda, serta menggunakan tanda
kurung secara berbeda (atau bahkan tidak menggunakannya sama sekali, seperti
dalam apa yang disebut notasi Polandia). Namun, semua pendekatan mempunyai
kesamaan bahwa kaidah pembentukan mendefinisikan himpunan suku dan !!{formula}
secara \emph{induktif}. Jika dilakukan dengan tepat, setiap ungkapan pada
dasarnya hanya dapat dihasilkan melalui satu cara menurut kaidah pembentukan.
Definisi induktif yang menghasilkan ungkapan yang \emph{dapat dibaca secara
unik} memungkinkan kita memberi makna kepada ungkapan-ungkapan tersebut
dengan metode yang sama---definisi induktif.

Pemberian makna kepada ungkapan merupakan ranah semantik. Konsep utama dalam
semantik ialah kepuasan dalam !!a{structure}. !!^a{structure} memberi makna
kepada unsur-unsur dasar bahasa: !!a{domain} ialah himpunan objek yang tak
kosong. Kuantor ditafsirkan merentang atas domain ini, !!{constant} diberi
elemen dalam domain, !!{function} diberi fungsi dari !!{domain} ke dirinya
sendiri, dan !!{predicate} diberi relasi pada !!{domain}. !!{domain} bersama
penetapan bagi kosakata dasar membentuk !!a{structure}. !!^{variable} dapat
muncul dalam !!{formula}; untuk memberikan semantiknya, kita juga harus
menetapkan !!{element} dari !!{domain} kepadanya---inilah penetapan variabel.
Terakhir, relasi kepuasan memadukan semua hal tersebut. !!^a{formula} dapat
terpenuhi dalam !!a{structure}~$\Struct{M}$ relatif terhadap penetapan
variabel~$s$, yang ditulis $\Sat{M}{!A}[s]$. Relasi ini juga didefinisikan
dengan induksi pada struktur~$!A$, menggunakan tabel kebenaran bagi penghubung
logis untuk mendefinisikan, misalnya, kepuasan $!A \land !B$ berdasarkan
terpenuhi atau tidak terpenuhinya $!A$ dan~$!B$. Ternyata penetapan variabel
tidak relevan jika !!{formula}~$!A$ merupakan !!a{sentence}, yakni tidak
mempunyai variabel bebas; karena itu kita dapat membicarakan !!{sentence} yang
terpenuhi (atau tidak) begitu saja dalam !!{structure}.

Berdasarkan relasi kepuasan $\Sat{M}{!A}$ bagi kalimat, selanjutnya kita dapat
mendefinisikan pengertian semantik dasar berupa validitas, konsekuensi, dan
keterpenuhan. Suatu kalimat valid, $\Entails !A$, jika setiap struktur
memenuhinya. Kalimat itu merupakan konsekuensi dari suatu himpunan
!!{sentence}, $\Gamma \Entails !A$, jika setiap !!{structure} yang memenuhi
semua !!{sentence} dalam~$\Gamma$ juga memenuhi~$!A$. Suatu himpunan kalimat
dapat dipenuhi jika terdapat !!{structure} yang memenuhi semua !!{sentence}
di dalamnya secara bersamaan. Karena !!{formula} didefinisikan secara
induktif, dan kepuasan pada gilirannya didefinisikan dengan induksi pada
struktur !!{formula}, kita dapat menggunakan induksi untuk membuktikan
sifat-sifat semantik kita dan menghubungkan pengertian-pengertian semantik
yang telah didefinisikan.

\end{document}
